\documentclass[11pt,a4paper]{article}

\usepackage[utf8]{inputenc}
\usepackage[T1]{fontenc}
\usepackage[brazilian]{babel}
\usepackage[margin=2cm]{geometry}
\usepackage{enumitem}
\usepackage{hyperref}
\usepackage{titlesec}
\usepackage{parskip}

\hypersetup{
  colorlinks=true,
  linkcolor=black,
  urlcolor=black,
  pdftitle={Kamille Hartmann Maccagnan - Currículo},
  pdfauthor={Kamille Hartmann Maccagnan}
}

\titleformat{\section}{\large\bfseries}{}{0em}{}[\titlerule]
\titlespacing*{\section}{0pt}{12pt}{6pt}

\pagestyle{empty}
\setlength{\parindent}{0pt}
\setlength{\parskip}{6pt}

\newcommand{\cvheader}[4]{%
  \begin{center}
    {\LARGE\bfseries #1}\\[4pt]
    {\large #2}\\[6pt]
    #3\\[2pt]
    #4
  \end{center}
  \vspace{4pt}
}

\newcommand{\experience}[3]{%
  \textbf{#1} \hfill \textit{#2}\\[4pt]
  #3
  \vspace{8pt}
}

\begin{document}

\cvheader{Kamille Hartmann Maccagnan}%
{Analista de BI Jr | Power BI | SQL | KPIs | Governança de Dados}%
{Campo Comprido, Curitiba-PR \textbullet{} (41) 9 9928-4424 \textbullet{} \href{mailto:kamillehartmannm@gmail.com}{kamillehartmannm@gmail.com}}%
{\href{https://linkedin.com/in/kamille-hartmann-maccagnan/}{linkedin.com/in/kamille-hartmann-maccagnan}}

\section*{Resumo Profissional}

Profissional de Business Intelligence com experiência no desenvolvimento e manutenção de dashboards, relatórios gerenciais e indicadores em Microsoft Power BI, apoiando a transformação de dados em informações para tomada de decisão. Uso diário de SQL para consultas, extração e consolidação de dados de diferentes fontes, com modelagem de informações, construção de métricas e KPIs em DAX e Power Query. Atuação no levantamento de requisitos junto às áreas de negócio, garantia de qualidade e consistência das informações, documentação técnica e funcional dos relatóios e suporte aos usuários na utilização das soluções analíticas. Vivência em governança de dados, identificação de inconsistências, melhoria contínua de processos de BI e interface com stakeholders em ambiente corporativo e multinacional. Uso de ferramentas de Inteligência Artificial, como ChatGPT, Copilot e recursos de insights e Q\&A do Power BI, como apoio à organização de ideias, análise e produtividade. Perfil analítico, organizado e com facilidade para traduzir necessidades de negócio em soluções de informação.

\section*{Experiência Profissional}

\experience{Anjuss -- Analista de Suporte de TI}{07/2025 -- Atual}{%
Atuação no desenvolvimento e manutenção de dashboards e painéis gerenciais em Power BI para monitoramento de indicadores operacionais e desempenho de processos, com construção de métricas, relatórios recorrentes e apoio à geração de insights para tomada de decisão. Uso diário de SQL para extração, consulta e consolidação de dados, complementado por transformações com Power Query para preparação e padronização de informações. Levantamento de necessidades junto às áreas de negócio e operação, com tradução de demandas em requisitos de informação e documentação funcional dos relatórios e indicadores desenvolvidos. Responsável pela identificação de inconsistências em bases de dados, garantia de qualidade e atualização das informações disponibilizadas aos usuários e prestação de suporte na utilização das soluções analíticas. Participação em iniciativas de melhoria contínua, automação de rotinas de BI e capacitação de usuários sobre indicadores e ferramentas. Utilização de ferramentas de Inteligência Artificial, como ChatGPT e Copilot, como apoio na elaboração de documentação e organização de análises.}

\experience{HORSE do Brasil -- Estagiária de TI (IS/IT LATAM)}{07/2024 -- 07/2025}{%
Atuação em governança de dados e Business Intelligence em ambiente multinacional, com elaboração de relatórios executivos e dashboards para monitoramento de indicadores de desempenho de projetos e iniciativas estratégicas. Consolidação de dados provenientes de diferentes fontes para apresentação clara a lideranças e stakeholders de negócio, fornecedores e times técnicos. Interface entre áreas de negócio e tecnologia para alinhamento de demandas analíticas, validação de entregas e compartilhamento de análises, com participação em reuniões conduzidas em até três idiomas simultaneamente. Atuação na documentação de processos, análise de riscos, padronização de informações e evolução de práticas de governança e qualidade de dados com foco em confiabilidade analítica.}

\experience{Restaurante Familiar -- Analista de Dados e Suporte Técnico}{01/2023 -- 07/2024}{%
Atuação no desenvolvimento de soluções analíticas para apoio à gestão do negócio, com extração, transformação e consolidação de dados de vendas, operação e resultados financeiros. Criação da base de dados do zero em Excel e modelagem de informações para viabilizar dashboards em Power BI, com definição de KPIs e indicadores de desempenho. Responsável pela padronização de dados, identificação de inconsistências, documentação de procedimentos e suporte às áreas na interpretação de relatórios e indicadores para tomada de decisão. Atuação no alinhamento entre áreas operacionais e administrativas, traduzindo necessidades de negócio em métricas e rotinas práticas de validação.}

\section*{Projetos e Conquistas}

Desenvolvimento de dashboards e painéis gerenciais para monitoramento de KPIs e suporte à tomada de decisão. Estruturação de bases de dados e automação de fluxos analíticos em Excel e Power BI, aumentando confiabilidade e reduzindo retrabalho. Participação em iniciativas de melhoria contínua e governança de dados. 1º lugar no Transformation Day Renault 2023 com solução orientada a dados e melhoria de processos. Participação no Hackathon Bosch 2023 com foco em resolução de problemas de negócio com uso de dados.

\section*{Habilidades Técnicas}

\textbf{Business Intelligence e Power BI:} Desenvolvimento e manutenção de dashboards e painéis gerenciais, modelagem de dados, DAX, Power Query, insights, Q\&A, relatórios gerenciais, construção de KPIs e métricas.

\textbf{SQL e Dados:} SQL (consultas, extração e consolidação diária), Excel avançado, extração, transformação e consolidação de dados (ETL com Power Query), noções de Data Warehouse e Business Intelligence.

\textbf{Governança e Qualidade:} Governança de dados, qualidade e consistência de informações, identificação de inconsistências, documentação técnica e funcional, melhoria contínua de processos de BI.

\textbf{Negócio e Requisitos:} Levantamento de requisitos com áreas de negócio, interpretação de dados, geração de insights, suporte a usuários, capacitação sobre soluções analíticas.

\textbf{Ferramentas:} Spotfire, Monday, Trello, Notion, ServiceNow, Pacote Office.

\textbf{Inteligência Artificial:} ChatGPT, Copilot, apoio à produtividade, organização de informações e análises.

\textbf{Metodologias:} Scrum, Kanban, Boas Práticas PMOBOK.

\textbf{Idiomas:} Inglês avançado \textbullet{} Espanhol básico \textbullet{} Coreano básico.

\section*{Capacitações Complementares}

Capacitação em aproveitamento de ferramentas de Inteligência Artificial.

\section*{Formação Acadêmica}

\textbf{PUCPR} -- Graduação em Gestão da Tecnologia da Informação \hfill Previsão de conclusão: 06/2026\\
Curitiba, PR

\textbf{Escola Conquer} -- Pós-graduação em Liderança e Gestão em Tecnologia \hfill Conclusão: 2024\\
Curitiba, PR

\textbf{Universidade Tuiuti do Paraná} -- Graduação em Medicina Veterinária \hfill Conclusão: 06/2019\\
Curitiba, PR

\textbf{Qualittas} -- Pós-graduação em Clínica Médica e Cirúrgica de Animais Exóticos e Pets Não Convencionais \hfill Conclusão: 12/2022\\
Curitiba, PR

\end{document}
